% Section 2.4: Symmetric and antisymmetric tensors; contraction of tensors | Part: 1 | Depends on: sec:vectors | Status: done | Original lines: 705-741
%!TEX root = ../../main.tex
\subsection{Symmetric and antisymmetric tensors; contraction of tensors}\label{sec:symcontr}
% REWRITTEN: the wedge product is defined once, in Section sec:ext (eq:wedge11); here it is only previewed.
A tensor $\tau\in\TT{0}{2}$ is \emph{symmetric} if $\tau(v,w)=\tau(w,v)$ and \emph{antisymmetric} if $\tau(v,w)=-\tau(w,v)$.
For covectors $\alpha,\beta$ the antisymmetric combination
\[
\alpha\otimes\beta-\beta\otimes\alpha
\]
will be called the \emph{wedge product} $\alpha\wedge\beta$ in Section~\ref{sec:ext}, where it is evaluated as
\eqref{eq:wedge11}: $(\alpha\wedge\beta)(v,w)=\alpha(v)\beta(w)-\alpha(w)\beta(v)$. Intuitively, this is the signed area of the
parallelogram spanned by $v,w$ as seen through the ``rulers'' $\alpha,\beta$. The notes use it in the area form of the sphere,
$r^2\sin\theta\,\dd\theta\wedge\dd\varphi$ (Example~\ref{ex:sphereareaform}).

\subhead{Contraction}
\begin{formal}[{contraction}]
The evaluation $V\times V^*\to\R$, $(v,f)\mapsto f(v)$ is bilinear, with $e_i,e^{j}\mapsto\delta^j{}_i$. Applying it to one covariant and one contravariant slot of a tensor gives a new tensor of type $(\ell-1,k-1)$. In components: set an upper and a lower index equal and sum,
\[
(c\tau)^{i_2\dots}{}_{j_2\dots}=\tau^{s\,i_2\dots}{}_{s\,j_2\dots}.
\]
This generalises the pairing \eqref{eq:pairing} of Section~\ref{sec:contraction}, which is the case $\tau=\omega\otimes v$.
\end{formal}
\begin{intuition}[{contraction = plug a vector-slot into a covector-slot}]
Think of connecting an output cable of one slot to the input of another and ``summing over all possible signals''. The result is lower in rank and \emph{still a tensor}. Examples: the trace $A^i{}_i$ of a linear map; the pairing $f_iv^i$; later the Ricci tensor $R^i{}_{ijk}$ and the scalar curvature.
\end{intuition}

% REORDERED: the exercises of this section (original L725-731) now close the section, in exercises.tex.

\begin{remark}[Three kinds of $2$-tensors]
A tensor in $\TT{0}{2}$ may be symmetric, antisymmetric, or neither; every one splits uniquely into a symmetric and an
antisymmetric part. A \emph{metric} is symmetric and non-degenerate (Section~\ref{sec:metric}); a \emph{$2$-form} is
antisymmetric (Section~\ref{sec:ext}); a general bilinear form such as $\dd x\otimes\dd y$ has both parts. The two parts
do different jobs -- one measures length, the other oriented area -- and the symmetric/antisymmetric bookkeeping
is made precise in Section~\ref{sec:wedgetensor}.
\end{remark}
